\documentclass[conference]{IEEEtran}

% ── Packages ─────────────────────────────────────────────────────────────────
\usepackage[T1]{fontenc}
\usepackage[utf8]{inputenc}
\usepackage{cite}
\usepackage{amsmath,amssymb}
\usepackage{graphicx}
\usepackage{url}
\usepackage[hidelinks]{hyperref}
\usepackage{booktabs}
\usepackage{array}
\usepackage{listings}
\usepackage{xcolor}
\usepackage{float}

% ── Code listing style ────────────────────────────────────────────────────────
\lstset{
  basicstyle=\ttfamily\footnotesize,
  keywordstyle=\color{blue}\bfseries,
  commentstyle=\color{gray},
  breaklines=true,
  frame=single,
  numbers=left,
  numberstyle=\tiny\color{gray},
  language=C
}

% ── Title ─────────────────────────────────────────────────────────────────────
\title{Concurrent Web Server: Thread Pool Architecture\\
with Size-Based Scheduling}

\author{%
  \begin{tabular}{cc}
    \begin{minipage}[t]{0.60\columnwidth}\centering
      \textbf{Ishit Arhatia}\\
      \small MS Data Science\\
      \small Northeastern University
    \end{minipage}
    &
    \begin{minipage}[t]{0.60\columnwidth}\centering
      \textbf{Rudrajit Banerjee}\\
      \small MS Computer Science\\
      \small Northeastern University
    \end{minipage}
    \\[2em]
    \begin{minipage}[t]{0.60\columnwidth}\centering
      \textbf{Helena Zhuo}\\
      \small MS Computer Science\\
      \small Northeastern University
    \end{minipage}
    &
    \begin{minipage}[t]{0.60\columnwidth}\centering
      \textbf{Tomas Colorado}\\
      \small MS Computer Science\\
      \small Northeastern University
    \end{minipage}
  \end{tabular}
}

\begin{document}

\maketitle

% ═════════════════════════════════════════════════════════════════════════════
\begin{abstract}
% ═════════════════════════════════════════════════════════════════════════════

This paper presents the design, implementation, and experimental evaluation
of a concurrent web server built around a fixed-size thread pool and a
producer-consumer bounded buffer, following the G.5 OSTEP specification. We
implement and compare two scheduling policies within the bounded buffer:
First-Come-First-Served (FCFS) as the baseline and Shortest-File-First (SFF)
as our research extension. Drawing on established theory that web file sizes
follow heavy-tailed distributions~\cite{harchol2003}, we hypothesized that
SFF would reduce mean response time without materially penalizing large-file
clients. Our measurements show that SFF provides modest improvements at low
concurrency ($+12\%$ for large-file, $+6.6\%$ for heavy-tailed at $c{=}10$)
but regresses at high concurrency ($-7.7\%$ to $-15.5\%$ at $c{=}100$),
attributable to the overhead of SFF's $O(n)$ buffer scan under mutex
contention at high load. Full results and analysis are presented in
Sections~\ref{sec:results} and~\ref{sec:analysis}.

\end{abstract}

\begin{IEEEkeywords}
Concurrent web server, thread pool, bounded buffer, producer-consumer,
scheduling policy, First-Come-First-Served (FCFS), Shortest-File-First (SFF),
size-based scheduling, heavy-tailed workloads, POSIX threads.
\end{IEEEkeywords}

% ═════════════════════════════════════════════════════════════════════════════
\section{Introduction}
% ═════════════════════════════════════════════════════════════════════════════

Modern web servers must handle thousands of simultaneous client connections
while maintaining low latency and high throughput. A foundational design choice
in any such system is how to move work from the single point of connection
acceptanc is a listening socket to a pool of threads that can service
requests in parallel. The canonical solution, described by Welsh et al.\ in the
SEDA architecture~\cite{welsh2001} and formalized in the OSTEP curriculum,
is the \emph{producer-consumer bounded buffer}: the main thread produces
connection descriptors into a fixed-capacity queue, and a fixed-size pool of
worker threads consume from it. This decoupling keeps the acceptance loop
fast and provides natural backpressure when workers are saturated.

Once this architecture is in place, a second design choice emerges: in what
order should workers dequeue requests? The simplest policy is FCFS
(First-Come-First-Served), which processes requests strictly in arrival order.
FCFS is fair, simple to implement, and free of starvation. However, web
workloads are not uniform. Crovella and Bestavros showed that World Wide Web
traffic is self-similar and that file sizes follow heavy-tailed (Pareto-like)
distributions~\cite{crovella1997}: the vast majority of objects are small
(a few kilobytes), yet a small number of large files account for most bytes
transferred. Under these conditions, FCFS performs poorly, a single large
file can block dozens of waiting small requests behind it.

Size-based scheduling policies such as Shortest-File-First (SFF) exploit this
skew. By allowing workers to select the smallest pending request from the
buffer rather than the oldest, small requests jump ahead of large ones,
completing quickly and clearing queue capacity. Harchol-Balter et al.\
proved theoretically and through trace-driven simulation that this approach
can reduce mean response time by orders of magnitude under moderate-to-high
load without materially starving large requests~\cite{harchol2003}. Schroeder
and Harchol-Balter further showed that the performance gap widens dramatically
as server load approaches capacity~\cite{schroeder2006}---exactly the operating
regime that matters most.

This paper makes the following contributions: (1) a fully conformant G.5
OSTEP concurrent web server with configurable scheduling policy; (2) an
automated benchmarking rig producing reproducible latency and throughput
measurements across three workload types and six concurrency levels; and
(3) an empirical comparison of FCFS and SFF scheduling under realistic
heavy-tailed workloads, evaluating whether the theoretical predictions of
Harchol-Balter hold in a practical POSIX thread pool implementation.

The remainder of this paper is organized as follows.
Section~\ref{sec:related} surveys the relevant literature.
Section~\ref{sec:design} describes the system architecture.
Section~\ref{sec:methodology} defines the experimental methodology.
Section~\ref{sec:results} presents the comparative results.
Section~\ref{sec:analysis} analyzes the findings.
Section~\ref{sec:conclusion} concludes and outlines future work.

% ═════════════════════════════════════════════════════════════════════════════
\section{Related Work}
\label{sec:related}
% ═════════════════════════════════════════════════════════════════════════════

Related work naturally divides into three themes: the architecture of
concurrent servers, the theory of size-based scheduling, and modern
production-scale scheduling systems that validate or extend these ideas.

\subsection{Architecture of Concurrent Servers}

The foundation of our design lies in the threading model introduced by
Anderson et al.~\cite{anderson1992}, which identified the core tension in
concurrent systems: kernel-level threads are expensive (context switches on
the order of 100\,$\mu$s) but integrate correctly with blocking I/O, while
user-level threads are an order of magnitude faster but stall the entire
process on any blocking system call. Our implementation resolves this tension
using POSIX \texttt{pthreads} with a fixed-size pool, accepting kernel-thread
cost in exchange for correct semantics when workers block on file I/O.

Building on this threading foundation, Welsh, Culler, and Brewer introduced
SEDA~\cite{welsh2001}; a Staged Event-Driven Architecture that partitions
server logic into a pipeline of stages, each with its own thread pool and
queue. SEDA's Haboob HTTP server used a form of shortest connection first
scheduling within each stage, directly motivating our SFF policy. The
producer-consumer bounded buffer at the heart of our design is a single-stage
simplification of SEDA's multi-stage pipeline.

von Behren et al.\ later challenged SEDA's event-driven philosophy with
Capriccio~\cite{vonbehren2003}, arguing that user-level thread systems can
match event-driven performance while retaining the simpler sequential
programming model. Capriccio introduced resource-aware scheduling, where the
scheduler used compiler-annotated resource usage predictions to make smarter thread selection decisions that are conceptually analogous to our use of
\texttt{stat()} to obtain file size before enqueuing a request.

\subsection{Size-Based Scheduling Theory}

The theoretical backbone of our research extension is the work of
Harchol-Balter et al.~\cite{harchol2003}, who formalized size-based
scheduling for web servers as an application of
Shortest-Remaining-Processing-Time (SRPT) queuing theory. The key insight is
that, under heavy-tailed workloads, prioritizing short jobs (small files)
dramatically reduces mean sojourn time because the population of competing
small jobs is large. The paper's trace-driven simulations showed mean response
time improvements of 2--7$\times$ over FCFS at moderate-to-high loads.

Schroeder and Harchol-Balter~\cite{schroeder2006} extended this analysis to
overload conditions, proving that FCFS response times explode exponentially as
load approaches server capacity, while SFF-style policies remain stable. The
paper also characterized the ``starvation myth'' wherein large requests are rarely
penalized significantly under SFF because they are statistically infrequent in heavy-tailed workloads.

The heavy-tailed nature of web traffic itself was characterized by Crovella
and Bestavros~\cite{crovella1997}, who demonstrated through empirical analysis
of HTTP traces that file size distributions are Pareto-like with shape
parameter $\alpha \approx 1.1$. This self-similarity is the fundamental
property that makes size-based scheduling effective and is the reason our
benchmarking rig generates workloads with mixed small and large files rather
than uniform sizes.

\subsection{Modern Production Scheduling Systems}

Recent systems work validates and extends the scheduling ideas above at
production scale. Ousterhout et al.'s Shenango~\cite{ousterhout2019} achieves
microsecond-scale thread scheduling via a dedicated IOKernel that reallocates
CPU cores across applications as request rates change. Its evaluation confirmed
that scheduling policy, not mechanism overhead, and is the primary determinant of tail latency, directly motivating our policy comparison.

Humphries et al.'s ghOSt~\cite{ghost2021}, deployed at Google, decouples
scheduling policy from kernel enforcement, allowing user-space agents to
implement FIFO, priority, or preemptive policies via kernel transactions. This
validates our design choice of making scheduling configurable via a runtime
flag rather than hardcoding one policy. Demoulin et al.'s
Perséphone~\cite{persephone2021} demonstrated that size-based scheduling
continues to provide substantial tail latency improvements for heavy-tailed
workloads at microsecond scale, corroborating the theoretical predictions our
experiment tests. Lampson's design heuristics~\cite{lampson1983} also informed
our architectural choices: splitting resources in a fixed way justifies our
fixed-size thread pool, and handling the normal and worst case separately
underpins our SFF rationale.

% ═════════════════════════════════════════════════════════════════════════════
\section{System Design}
\label{sec:design}
% ═════════════════════════════════════════════════════════════════════════════

\subsection{Architectural Overview}

Our system implements the G.5 OSTEP concurrent web server specification. The
server accepts connections on a TCP socket and distributes HTTP requests to a
fixed pool of worker threads via a shared bounded buffer, following the
producer-consumer pattern. Three configurable parameters govern runtime
behavior: \texttt{-t~N} sets the thread pool size (default: 4), \texttt{-b~N}
sets the buffer capacity (default: 16), and \texttt{-s~\{FCFS|SFF\}} selects
the scheduling policy.

\subsection{Thread Pool and Bounded Buffer}

The bounded buffer is implemented as a fixed-capacity circular array protected
by a single \texttt{pthread\_mutex\_t} and two condition variables:
\texttt{not\_full} (signaled by consumers after dequeue) and \texttt{not\_empty}
(signaled by the producer after enqueue). A fixed-size array is used rather
than a dynamic queue to ensure memory-safe operation under concurrent access
without requiring heap allocation inside the critical section.

Each slot in the buffer holds a \texttt{request\_t} struct that carries all
the information a worker thread needs to service a request without accessing
any shared state after dequeue. The \texttt{client\_fd} field is the accepted
socket descriptor; \texttt{filename} holds the sanitized URI path;
\texttt{file\_size} stores the result of the \texttt{stat()} call made before
acquiring the mutex; and \texttt{arrival\_ms} records the enqueue timestamp
for FCFS ordering. The struct is defined as follows:

\begin{lstlisting}
typedef struct {
    int    client_fd;     /* accepted socket descriptor */
    char   filename[256]; /* parsed URI path           */
    off_t  file_size;     /* from stat(), set on enqueue */
    long   arrival_ms;    /* timestamp for FCFS order  */
} request_t;
\end{lstlisting}

\begin{figure}[t]
  \centering
  \includegraphics[width=\columnwidth]{figs/bounded-buffer-activity.pdf}
  \caption{Activity diagram of the bounded buffer. The listener thread
    (producer) enqueues connection descriptors; worker threads (consumers)
    dequeue according to the active scheduling policy.}
  \label{fig:activity}
\end{figure}

\subsection{Scheduling Policies}

\textbf{FCFS.} The consumer removes the element at the buffer's logical
head which is a standard circular-queue dequeue in $O(1)$.

\textbf{SFF.} Rather than scanning at dequeue time, SFF maintains a sorted
buffer by reorganizing at \emph{insert} time. When the producer enqueues a
new request, it acquires the lock, then scans backwards from the last occupied
slot toward the head, shifting entries with a larger \texttt{file\_size} one
slot forward until it finds an entry less than or equal to the new request (or
exhausts all entries). The new request is inserted at that position, keeping
the buffer sorted in ascending \texttt{file\_size} order at all times. As a
result, \texttt{buffer\_get} always takes from the head in $O(1)$, the
shortest file is always first. The scheduling selection is encapsulated such
that the worker code is entirely policy-agnostic; only the producer's insert
path differs between FCFS and SFF.

\subsection{HTTP Request--Response Lifecycle}

\begin{figure}[t]
  \centering
  \includegraphics[width=\columnwidth]{figs/http-sequence.pdf}
  \caption{Sequence diagram of the HTTP request--response lifecycle across
    five steps: TCP handshake, request transmission, server processing,
    response delivery, and connection teardown.}
  \label{fig:sequence}
\end{figure}

On connection establishment, the main thread completes the TCP three-way
handshake and enqueues the connected descriptor into the bounded buffer. A
worker thread dequeues it, parses the HTTP request line; under SFF the
\texttt{stat()} call is made by the producer at enqueue time to obtain
\texttt{file\_size} for ordering, while under FCFS the worker performs
\texttt{stat()} after dequeue to verify the file and obtain its size. The
worker then transmits the file with an appropriate status line. The server returns \texttt{200~OK} for
valid requests and \texttt{404~Not Found} for missing resources.

\subsection{Request-to-Job Dispatch}

\begin{figure}[t]
  \centering
  \includegraphics[width=\columnwidth]{figs/request-job-lifecycle.pdf}
  \caption{Request-to-job lifecycle. The research intervention occurs at
    Step~4 (Job Dispatch), where the Dispatcher applies FCFS or SFF.}
  \label{fig:lifecycle}
\end{figure}

Figure~\ref{fig:lifecycle} illustrates the system from a kernel-interaction
perspective. The Dispatcher is the bounded buffer's dequeue logic is that the precise point where scheduling policy determines which request is served next.

\subsection{Source Organization}

\begin{table}[h]
\centering
\caption{Source file responsibilities}
\label{tab:files}
\small
\begin{tabular}{@{}ll@{}}
\toprule
\textbf{File} & \textbf{Responsibility} \\
\midrule
\texttt{main.c}              & Arg parsing, socket setup, accept loop \\
\texttt{bounded\_buffer.c/h} & Buffer, mutex, condvars, FCFS/SFF \\
\texttt{thread\_pool.c/h}    & Pool init, thread creation, teardown \\
\texttt{http.c/h}            & HTTP parsing, file send, status codes \\
\texttt{mock\_client.c}      & Correctness testing client \\
\texttt{benchmark.sh}        & Automated load-testing script \\
\texttt{Makefile}            & \texttt{make} / \texttt{make clean} \\
\bottomrule
\end{tabular}
\end{table}

\subsection{Concurrency Correctness}

All accesses to the shared bounded buffer pass through the mutex. Condition
variables are checked inside \texttt{while} loops to guard against spurious
wakeups. The producer holds the mutex only while inserting into the array;
the \texttt{stat()} call is performed beforehand. Workers hold the mutex only
while scanning/removing from the buffer; all file I/O occurs after releasing
the lock, minimizing critical section duration and avoiding holding the mutex
across blocking calls.

% ═════════════════════════════════════════════════════════════════════════════
\section{Methodology}
\label{sec:methodology}
% ═════════════════════════════════════════════════════════════════════════════

\subsection{Benchmarking Rig}

All performance measurements are collected externally, treating the server
as a black box. We use \texttt{wrk}~\cite{wrk} as our primary load generator
because it supports configurable concurrency, duration-based runs, and Lua
scripting for mixed-file-size workloads. Scripts are located in the project
repository~\cite{repo} under \texttt{/benchmarks/} and are fully automated:
a single invocation of \texttt{benchmark.sh} creates the small and large test
files using the \texttt{dd} Unix utility, starts the server, runs all
experiments, tears down the server, and writes raw output to
\texttt{/benchmarks/raw/}.

\subsection{Experimental Variables}

\textbf{Independent variables:}
\begin{itemize}
  \item \textbf{Scheduling policy:} FCFS vs.\ SFF (primary comparison).
  \item \textbf{Concurrency level:} $\{1, 10, 25, 50, 100, 200\}$ simultaneous
    \texttt{wrk} connections.
  \item \textbf{File size mix:} uniform small (4\,KB), uniform large (1\,MB),
    and heavy-tailed (90\% small / 10\% large by request count).
\end{itemize}

\textbf{Dependent variables:} throughput (req/s); response time mean, p50,
p95, p99 (ms); error rate (connection resets, timeouts).

\textbf{Controlled variables:} thread pool size fixed at 4; buffer capacity
fixed at 16; same VM and OS image; no background processes during measurement;
server and client run on the same host to eliminate network variability.

\subsection{Experimental Protocol}

Each workload type is run independently at each concurrency level. Each
configuration runs for 30 seconds with a per-trial warm-up period of 5
seconds (excluded from measurement). The server is restarted between trials
so each trial starts cold and requires its own warm-up. We collect five
independent trials per configuration and report the mean. Configurations
where the server reaches connection saturation are noted as such rather than
omitted, as saturation behavior is itself a meaningful data point for
comparing FCFS and SFF under overload.

\subsection{AI-Assisted Development}

Our team used AI assistance (Claude, Anthropic) throughout the project in
three distinct ways. First, for paper drafting and iteration: the IEEE
LaTeX structure, section prose, and bibliography were co-developed with AI,
with the team reviewing and correcting all factual claims and citations.
Second, for debugging: AI was consulted when diagnosing the server's
connection-refusal behavior under high concurrency, suggesting \texttt{ulimit}
adjustments and fd leak analysis, though the root cause was ultimately
identified through manual \texttt{htop} inspection. Third, for benchmark
interpretation: AI assisted in averaging raw CSV trial data and identifying
unit inconsistencies (microseconds vs.\ milliseconds) in \texttt{wrk} output.
The primary limitation observed was AI's tendency to hallucinate
implementation-level details that did not match the actual codebase;
all code was written and verified by team members directly.

\subsection{Environment}

\begin{table}[h]
\centering
\caption{Measurement environment}
\label{tab:env}
\begin{tabular}{@{}ll@{}}
\toprule
\textbf{Component} & \textbf{Specification} \\
\midrule
Host machine    & Windows 10 \\
Virtualization  & Oracle VirtualBox \\
Guest OS        & Ubuntu 22.04 LTS (Linux 5.15) \\
Guest resources & 4 vCPUs, 8\,GB RAM \\
Compiler        & GCC 11.4, \texttt{-O2} \\
Load generator  & \texttt{wrk} 4.1.0 \\
\bottomrule
\end{tabular}
\end{table}

% ═════════════════════════════════════════════════════════════════════════════
\section{Results}
\label{sec:results}
% ═════════════════════════════════════════════════════════════════════════════

\subsection{FCFS Baseline}

Table~\ref{tab:baseline} shows FCFS measurements under the three file-size
workloads at concurrency levels 10 and 100. Values represent the mean across
five independent trials. Raw \texttt{wrk} output is committed to
\texttt{/benchmarks/raw/}.

\begin{table}[h]
\centering
\caption{FCFS baseline: mean response time (ms) and throughput (req/s)}
\label{tab:baseline}
\begin{tabular}{@{}ccccc@{}}
\toprule
\textbf{Workload} & \textbf{Conc.} & \textbf{Mean RT} & \textbf{p99} & \textbf{Req/s} \\
\midrule
Small (4\,KB)  &  10 & 0.56\,ms  & 3.15\,ms  & 15{,}122 \\
Small (4\,KB)  & 100 & 3.55\,ms  & 6.81\,ms  & 26{,}715 \\
Large (1\,MB)  &  10 & 6.64\,ms  & 17.97\,ms & 1{,}462  \\
Large (1\,MB)  & 100 & 56.34\,ms & 87.61\,ms & 1{,}765  \\
Heavy-tailed   &  10 & 1.21\,ms  & 5.97\,ms  & 8{,}074  \\
Heavy-tailed   & 100 & 8.79\,ms  & 16.23\,ms & 11{,}149 \\
\bottomrule
\end{tabular}
\end{table}

\subsection{FCFS vs.\ SFF Comparison}

Table~\ref{tab:comparison} presents the side-by-side comparison of FCFS and
SFF across the same configurations. SFF results are from Sprint~4 benchmark
runs using identical methodology.

\begin{table}[h]
\centering
\caption{FCFS vs.\ SFF: mean response time (ms)}
\label{tab:comparison}
\begin{tabular}{@{}cccccc@{}}
\toprule
\textbf{Workload} & \textbf{Conc.} & \textbf{FCFS} & \textbf{SFF} & \textbf{Gain} \\
\midrule
Small (4\,KB)  &  10 & 0.56  & 0.56  & $\approx$0\%  \\
Small (4\,KB)  & 100 & 3.55  & 4.10  & $-15.5\%$     \\
Large (1\,MB)  &  10 & 6.64  & 5.84  & $+12.0\%$     \\
Large (1\,MB)  & 100 & 56.34 & 60.66 & $-7.7\%$      \\
Heavy-tailed   &  10 & 1.21  & 1.13  & $+6.6\%$      \\
Heavy-tailed   & 100 & 8.79  & 10.00 & $-13.8\%$     \\
\bottomrule
\end{tabular}
\end{table}

% ═════════════════════════════════════════════════════════════════════════════
\section{Analysis and Discussion}
\label{sec:analysis}
% ═════════════════════════════════════════════════════════════════════════════

\subsection{FCFS Baseline Behavior}

The FCFS baseline results confirm several theoretical predictions. Under the
small-file workload, mean response time scales modestly with concurrency
(0.56\,ms at $c{=}10$ to 3.55\,ms at $c{=}100$) while throughput increases
toward a plateau near 27{,}000\,req/s, consistent with saturation of the 4
worker threads. Under the large-file workload, the mean response time increase
from 6.64\,ms to 56.34\,ms (an $8.5\times$ jump) as concurrency increases
from 10 to 100 reflects each worker thread being tied up for the full duration
of a 1\,MB file transfer, causing the queue to build. The heavy-tailed workload
sits between these extremes, with mean response time rising from 1.21\,ms to
8.79\,ms which is evidence that the 90\% small-file majority keeps throughput
high, while the 10\% large-file minority begins to dominate latency at higher
concurrency.

\subsection{SFF vs.\ FCFS Comparison}

The results partially confirm our hypothesis but reveal an important
interaction between concurrency level and scheduling policy. At
concurrency~10, SFF modestly outperforms FCFS for large-file
($+12.0\%$ mean RT reduction) and heavy-tailed ($+6.6\%$) workloads,
consistent with the theoretical prediction that size-based scheduling
benefits emerge when a mix of small and large requests compete for the
same workers. For the uniform-small workload at concurrency~10, the two
policies perform identically, as expected, since all requests have the
same size and SFF degenerates to an arbitrary selection.

However, at concurrency~100, SFF \emph{underperforms} FCFS across all
three workload types: $-15.5\%$ for small files, $-7.7\%$ for large
files, and $-13.8\%$ for heavy-tailed. We attribute this regression to
two factors. First, with a buffer capacity of 16 slots and 4 worker
threads, the buffer fills quickly under 100 concurrent connections;
SFF's sorted-insert policy requires the producer to shift existing
buffer entries under the mutex, adding lock-hold time at enqueue that
delays competing producers. Under high concurrency this contention is
frequent, increasing effective enqueue latency. Second, at high
concurrency all four worker threads are almost always busy, meaning the
scheduling advantage of having the smallest file at the head is crowded
out by raw thread saturation --- workers dequeue and immediately process
regardless of ordering, so the insert-time cost is paid without the
full scheduling benefit being realized.

\subsection{Limitations}

Our SFF implementation performs an $O(n)$ sorted insert at enqueue time
rather than using a heap. For the buffer sizes used in this experiment
(16 slots), this is manageable, but it would become a bottleneck at
larger buffer capacities. Additionally, the \texttt{stat()} call is made
by the producer before enqueue under SFF, meaning file size is measured
at admission time; if files change between enqueue and dequeue, the size
estimate may be stale. Finally, running the server and the load generator
on the same VM host means CPU contention between \texttt{wrk} and the
server may inflate measured latencies.

% ═════════════════════════════════════════════════════════════════════════════
\section{Conclusion and Future Work}
\label{sec:conclusion}
% ═════════════════════════════════════════════════════════════════════════════

This paper presented the design, implementation, and empirical evaluation of
a concurrent POSIX web server with configurable FCFS and SFF scheduling. Our
system correctly serves static files to concurrent clients using a fixed-size
thread pool, a producer-consumer bounded buffer, and two scheduling policies
selectable at runtime. Three UML diagrams document the critical paths through
the system, and an automated benchmarking rig produces reproducible results
across six concurrency levels and three workload types.

The FCFS baseline confirms that large-file requests are the dominant driver
of latency degradation under load, with mean response time rising $8.5\times$
from concurrency~10 to 100. The heavy-tailed workload demonstrates that a
realistic web traffic mix sits between the two extremes, making it the most
interesting regime for the FCFS--SFF comparison.

The related work survey established that size-based scheduling is theoretically
well-motivated~\cite{harchol2003,schroeder2006}, architecturally
principled~\cite{welsh2001,lampson1983}, and validated at production
scale~\cite{ghost2021,persephone2021,ousterhout2019}.

\textbf{Future Work.} We would address three limitations identified above.
First, replacing the $O(n)$ sorted insert with a min-heap would make
SFF's enqueue operation $O(\log n)$ and reduce lock-hold time under high
concurrency. Second, implementing \texttt{keep-alive}
connection reuse would reduce TCP handshake overhead at high concurrency,
exposing the scheduling signal more clearly. Third, running the server and
load generator on separate machines would eliminate CPU contention and produce
cleaner latency measurements. A further extension would implement
Shortest-Remaining-Processing-Time (SRPT) rather than SFF, tracking remaining
bytes rather than total file size for even finer-grained prioritization.

% ═════════════════════════════════════════════════════════════════════════════
\begin{thebibliography}{99}
% ═════════════════════════════════════════════════════════════════════════════

\bibitem{repo}
I.~Arhatia, R.~Banerjee, H.~Zhuo, and T.~Colorado,
``project-webserver-team0101,''
GitHub, 2026.
\url{https://github.com/cs5600-sp26/project-webserver-team0101}

\bibitem{ostep}
R. H. Arpaci-Dusseau and A. C. Arpaci-Dusseau,
\textit{Operating Systems: Three Easy Pieces}, v1.10.
Arpaci-Dusseau Books, 2023.
\url{https://pages.cs.wisc.edu/~remzi/OSTEP/}

\bibitem{welsh2001}
M. Welsh, D. Culler, and E. Brewer,
``SEDA: An Architecture for Well-Conditioned, Scalable Internet Services,''
in \textit{Proc. 18th ACM Symp. Operating Systems Principles (SOSP~'01)},
Banff, Canada, Oct. 2001, pp.~230--243.

\bibitem{harchol2003}
M. Harchol-Balter, B. Schroeder, N. Bansal, and M. Agrawal,
``Size-Based Scheduling to Improve Web Performance,''
\textit{ACM Trans. Comput. Syst.}, vol.~21, no.~2, pp.~207--233, May 2003.

\bibitem{schroeder2006}
B. Schroeder and M. Harchol-Balter,
``Web Servers Under Overload: How Scheduling Can Help,''
\textit{ACM Trans. Internet Technol.}, vol.~6, no.~1, pp.~20--52, Feb. 2006.

\bibitem{crovella1997}
M. E. Crovella and A. Bestavros,
``Self-Similarity in World Wide Web Traffic: Evidence and Possible Causes,''
\textit{IEEE/ACM Trans. Networking}, vol.~5, no.~6, pp.~835--846, Dec. 1997.

\bibitem{anderson1992}
T. E. Anderson, B. N. Bershad, E. D. Lazowska, and H. M. Levy,
``Scheduler Activations: Effective Kernel Support for the User-Level
Management of Parallelism,''
\textit{ACM Trans. Comput. Syst.}, vol.~10, no.~1, pp.~53--79, Feb. 1992.

\bibitem{vonbehren2003}
R. von Behren, J. Condit, F. Zhou, G. C. Necula, and E. Brewer,
``Capriccio: Scalable Threads for Internet Services,''
in \textit{Proc. 19th ACM Symp. Operating Systems Principles (SOSP~'03)},
Bolton Landing, NY, Oct. 2003, pp.~268--281.

\bibitem{lampson1983}
B. W. Lampson,
``Hints for Computer System Design,''
\textit{ACM SIGOPS Oper. Syst. Rev.}, vol.~17, no.~5, pp.~33--48, Oct. 1983.

\bibitem{ghost2021}
J. T. Humphries \textit{et al.},
``ghOSt: Fast \& Flexible User-Space Delegation of Linux Scheduling,''
in \textit{Proc. 28th ACM Symp. Operating Systems Principles (SOSP~'21)},
Virtual, Oct. 2021, pp.~588--604.

\bibitem{persephone2021}
H. M. Demoulin \textit{et al.},
``When Idling Is Ideal: Optimizing Tail-Latency for Heavy-Tailed Datacenter
Workloads with Perséphone,''
in \textit{Proc. 28th ACM Symp. Operating Systems Principles (SOSP~'21)},
Virtual, Oct. 2021, pp.~621--637.

\bibitem{ousterhout2019}
A. Ousterhout, J. Fried, J. Behrens, A. Belay, and H. Balakrishnan,
``Shenango: Achieving High CPU Efficiency for Latency-sensitive Datacenter
Workloads,''
in \textit{Proc. 16th USENIX Symp. Networked Systems Design and
Implementation (NSDI~'19)},
Boston, MA, Feb. 2019, pp.~361--378.

\bibitem{wrk}
W. Glozer, ``wrk --- a HTTP benchmarking tool,'' GitHub, 2023.
\url{https://github.com/wg/wrk}

\end{thebibliography}

\end{document}